\begin{proposition}[P4.22]\hypertarget{tag-02WU}{}\label{night:P4.22}\label{stmt:P4.22}
\StatementRecord{night:P4.22}{proved}{02WU}{\SeriesRef{night:P4.13}, \SeriesRef{night:P4.15}, \SeriesRef{night:P4R.1}, \SeriesRef{night:P4R.3}}

In the standing torus-fibration setting of this section, an irreducible singular support
$S$ whose smooth normalization $T$ is birationally ruled
over an elliptic curve has $e(S)\ge0$. Equality holds
exactly when $T$ is minimal.
\end{proposition}
\begin{proof}

\begin{enumerate}[label={$\langle1\rangle$\arabic*.},leftmargin=*]
\item Push the effective conductor to the minimal ruled
surface $T_0\to B$, where $B$ is elliptic. We have
$K_{T_0}+D_0\equiv0$. If a horizontal component $H$
has coefficient one, adjunction gives $p_a(H)\le1$,
whereas its nonconstant map to $B$ gives geometric
genus at least one. Thus $H$ is smooth elliptic and
$(D_0-H)\cdot H=0$. There is no vertical component.
The horizontal part is one elliptic bisection or two
disjoint elliptic sections, all with coefficient one.
All subsequent blowups have smooth coefficient-one
conductor centres; their exceptional curves are absent
from the conductor. The normalization differential
has no divisorial rank drop. P4.15 makes $T$ minimal,
makes the differential immersive everywhere, and
gives $e(S)=0$.

\item The other possibility is
\[
 D_0=2H+\sum_j a_jF_j,
\]
where $H$ is a smooth elliptic section. As in P4R.3,
$D_{\mathrm{red}}$ is $H$ with rational trees attached,
and its inclusion in $T$ induces an isomorphism on
fundamental groups. The conductor pushout and P4R.1
make $\pi_1(C)$ virtually abelian. The reduced-curve
calculation in P4R.3 gives $e(C)\ge0$, and P4.13 then
gives $e(S)\ge0$.

\item If $e(S)=0$, the curve calculation forces $C$ to have
one component, genus-one smooth normalization, and
no point-identification graph cycles. Each conductor
component maps finitely and nonconstantly onto $C$.
Riemann--Hurwitz excludes rational conductor
components. Thus $D_{\mathrm{red}}=H$ and no blowup
occurred: the first blowup on $2H$ would create a
vertical conductor component of coefficient one.
Therefore $T$ is minimal.

\item Conversely suppose $T$ is minimal. In the double-section
case let $v$ be the number of distinct vertical fibres
in the conductor. Since $e(T)=e(H)=0$ and each vertical
$\mathbb P^1$ meets $H$ once,
\[
 e(S_{\mathrm{reg}})=e(T\setminus D_{\mathrm{red}})=-v.
\]
P4.13 forces $v=0$, hence $D=2H$. At a general point
of its image curve, every normalization branch has
conductor exponent two. The branch exponent equals
its intrinsic conductor exponent plus the sum of its
intersections with the other branches. A singular
branch together with another branch would have
exponent greater than two. Thus either there is one
ordinary cusp branch, or all branches are smooth:
two with intersection multiplicity two, or three
with pairwise intersection multiplicity one. Four
smooth branches would have exponent at least three,
so the list is exhaustive.

\item In the cusp case $H\to C$ has generic degree one.
The normalization of $C$ is therefore $H$, an elliptic
curve. Virtual abelianness of $\pi_1(C)$ excludes any
additional point-identification graph cycle, so
$e(C)=0$. Additivity gives $e(S)=0$. In the other two
cases the normalization differential is immersive at
the generic point of $H$. There is no other conductor
component, so its rank-drop set is finite. P4.15
applies and again gives $e(S)=0$.

\end{enumerate}
\end{proof}